\section{Related work}
\paragraph{GUI agents and datasets.} Large demonstration corpora such as AITW~\citep{aitw} and AndroidControl~\citep{androidcontrol} pair natural-language instructions with screen--action sequences. AITZ~\citep{aitz} and AMEX~\citep{amex} add screen descriptions, action rationales and element-function annotations, and GUI-Odyssey~\citep{odysseypub} collects longer cross-app tasks. AndroidWorld~\citep{androidworld} instead runs tasks on a live device and checks the final system state. On the modeling side, Auto-GUI~\citep{autogui} predicts the next action from the screen and the previous actions; CogAgent~\citep{cogagent}, SeeClick~\citep{seeclick} and OS-Atlas~\citep{osatlas} train on large-scale GUI grounding data to locate interface elements; and UI-TARS~\citep{uitars} maps screenshots to actions end to end. These works mainly address where to act on the current screen. We study a different kind of localization, in which the answer is no longer on the current screen and the agent must find the earlier screen that shows it.

\paragraph{History and memory in GUI agents.} Many agents keep the history as a textual action sequence~\citep{autogui} or compress past screenshots; GUI-Odyssey's OdysseyAgent, for example, resamples historical screenshot features~\citep{odysseypub}. Other agents record information while they act: Mobile-Agent-v2~\citep{mobileagentv2} has a decision agent update a memory unit of task-relevant content, Mobile-Agent-E~\citep{mobileagente} uses a notetaker agent to record important information, and GUI-Rise~\citep{guirise} writes a history summary at every step and trains it with a reward tied to later actions. All of them decide what to keep when a screen is first seen, before the agent knows which field it will need. Cross-episode memories, such as AppAgent's exploration documents~\citep{appagent}, Synapse's exemplar trajectories~\citep{synapse} and the reusable workflows of Agent Workflow Memory~\citep{awm}, store how to perform tasks rather than what the current episode showed. Closest to our work is PAL-UI~\citep{palui}, which summarizes the history and trains an agent to retrieve an earlier screenshot when needed. Its 8.6K step-level examples are synthesized by a teacher model and supervise both when to look back and which frame to retrieve. We isolate the second decision and supervise it with the text recorded in the demonstration, so any corpus that logs typed text can supply training data.

\paragraph{Long-context and multi-image retrieval.} Lost in the Middle~\citep{lost} finds that language models make poor use of information in the middle of long inputs. RULER~\citep{ruler} shows that simple needle-in-a-haystack tests overstate the context length that models can actually use, and Visual Haystacks~\citep{haystacks} extends the question to large image collections, finding that models are misled by visual distractors and show a positional bias. ReadAgent~\citep{readagent} and MemGPT~\citep{memgpt} work within a limited context by retrieving original content through gists or a hierarchical memory. Our histories span only 1--5K tokens, so a full-history prompt remains feasible. The answer, however, often recurs across screens, which makes look-back distance ambiguous; Section~\ref{sec:distance} therefore reports the distance to both the earliest and the most recent occurrence.

\paragraph{Distant supervision from known answers.} Distant supervision was introduced for relation extraction, where facts in a knowledge base label sentences automatically~\citep{mintz}. Open-domain question answering adopted the idea: DrQA~\citep{drqa} trains its reader on retrieved paragraphs that contain the answer string, ORQA~\citep{orqa} learns a retriever and a reader jointly from question--answer pairs alone, and DPR~\citep{dpr} uses BM25 passages that contain the answer as retriever positives when no gold passage is given. We bring this idea to GUI trajectories. The typed text plays the role of the answer and the screens before the current step bound the search, so we can check every earlier screen and keep all matches rather than only the top-ranked one. Our two selectors also mirror the two common reranking formulations: the per-screen scorer scores each candidate independently, as monoT5~\citep{monot5} does with the probabilities of \texttt{true} and \texttt{false}, while the full-history selector compares all candidates in one prompt, as RankGPT~\citep{rankgpt} does. A match, however, only shows that a screen contains the text; picking the right field on that screen remains a separate reading problem.

\section{Code and reproduction}
All code, derived data and per-case predictions are released in the \href{https://github.com/WMSlalalala/verifier-supervised-lookback}{project repository}. Its \texttt{src/} directory holds the scripts that sample GUI-Odyssey trajectories, run OCR and verify cases, build the case--screen labels (\texttt{build\_scorer\_data.py}), train the two selectors (\texttt{train\_scorer.py} and \texttt{train\_fullhist.py}) and run the reader (\texttt{read\_scorer.py}), together with the analysis scripts behind every table and figure. EasyOCR provides text recognition, and Transformers and PEFT run the Qwen2.5 models and their LoRA adapters.

Because the OCR records, labels and predictions are released in \texttt{data/} and \texttt{results/}, the main numbers can be recomputed on a CPU without the original screenshots or model weights. The snippet below uses only the Python standard library and prints 188/123/80, 188/136/84, 93/71/48 and 93/74/46 (cases, correct screens, correct typed text) for the trained full-history selector and the scorer on the development and held-out sets, matching Table~\ref{tab:main}.
\begin{lstlisting}[language=Python]
import json
from pathlib import Path
for split in ("", "_more"):
    for model in ("fullhist", "scorer"):
        path = Path("results") / f"{model}_trained{split}_hist_e2e.json"
        rows = json.loads(path.read_text(encoding="utf-8"))
        print(model, len(rows), sum(r["loc"] for r in rows),
              sum(r["e2e"] for r in rows))
\end{lstlisting}
Retraining both selectors over the five folds and running the reader needs a CUDA GPU (the reported 7B runs used a 48~GB card); for the held-out set, set \texttt{FOLD=ext} for the trainers and \texttt{SET=\_more} for the reader. The repository README lists the remaining commands and maps each table and figure in this report to the files that produce it.
\begin{lstlisting}[language=bash]
mkdir -p run-data && cp data/*.json* results/*.json run-data/
bash configure.sh "$PWD/run-data"
for f in 0 1 2 3 4; do
  FOLD=$f python src/train_scorer.py
  FOLD=$f python src/train_fullhist.py
done
for model in scorer fullhist; do
  PREFIX=$model SRC=trained HIST=1 python src/read_scorer.py
done
\end{lstlisting}

\section{Conclusion and limitations}
A recorded text entry already tells us which earlier screens hold the answer. Using these matches as labels, without human annotation or a teacher model, a 7B per-screen scorer finds the source screen in 72.3\% of development cases instead of 50.0\% and in 79.6\% of held-out cases instead of 41.9\%, and a full-history selector trained on the same labels performs comparably. Reading remains the bottleneck: the reader still misses the value in 62 of the scorer's 136 correct selections.

Our evidence has four main limitations. All four samples come from one corpus and one OCR engine. A substring can match by coincidence, and OCR errors or paraphrases can hide genuine matches. A 20-case spot check found the target visible on every inspected source screen, and in five of these cases the user had typed it earlier; this supports the presence of visible evidence but says little about whether users actually needed to remember the value. Finally, localization uses the same criterion that creates the labels, typed-text accuracy accepts the reference as a substring of a longer generation, and neither metric measures whole-task success.

\begingroup\small\setlength{\bibsep}{2pt}
\begin{thebibliography}{32}
\bibitem{appagent} C. Zhang, Z. Yang, J. Liu, Y. Han, X. Chen, Z. Huang, B. Fu and G. Yu. \href{https://arxiv.org/abs/2312.13771}{AppAgent: Multimodal Agents as Smartphone Users}. CHI, 2025.
\bibitem{mobileagentv2} J. Wang, H. Xu, H. Jia, X. Zhang, M. Yan, W. Shen, J. Zhang, F. Huang and J. Sang. \href{https://proceedings.neurips.cc/paper_files/paper/2024/hash/0520537ba799d375b8ff5523295c337a-Abstract-Conference.html}{Mobile-Agent-v2: Mobile Device Operation Assistant with Effective Navigation via Multi-Agent Collaboration}. NeurIPS, 2024.
\bibitem{uitars} Y. Qin et al. \href{https://arxiv.org/abs/2501.12326}{UI-TARS: Pioneering Automated GUI Interaction with Native Agents}. arXiv:2501.12326, 2025.
\bibitem{aitw} C. Rawles, A. Li, D. Rodriguez, O. Riva and T. Lillicrap. \href{https://arxiv.org/abs/2307.10088}{Android in the Wild: A Large-Scale Dataset for Android Device Control}. NeurIPS Datasets and Benchmarks Track, 2023.
\bibitem{androidcontrol} W. Li et al. \href{https://arxiv.org/abs/2406.03679}{On the Effects of Data Scale on UI Control Agents}. NeurIPS Datasets and Benchmarks Track, 2024.
\bibitem{odysseypub} Q. Lu, W. Shao, Z. Liu, L. Du, F. Meng, B. Li, B. Chen, S. Huang, K. Zhang and P. Luo. \href{https://openaccess.thecvf.com/content/ICCV2025/html/Lu_GUIOdyssey_A_Comprehensive_Dataset_for_Cross-App_GUI_Navigation_on_Mobile_ICCV_2025_paper.html}{GUIOdyssey: A Comprehensive Dataset for Cross-App GUI Navigation on Mobile Devices}. ICCV, 2025.
\bibitem{lost} N. F. Liu, K. Lin, J. Hewitt, A. Paranjape, M. Bevilacqua, F. Petroni and P. Liang. \href{https://aclanthology.org/2024.tacl-1.9/}{Lost in the Middle: How Language Models Use Long Contexts}. TACL, 12:157--173, 2024.
\bibitem{ruler} C.-P. Hsieh et al. \href{https://arxiv.org/abs/2404.06654}{RULER: What's the Real Context Size of Your Long-Context Language Models?} COLM, 2024.
\bibitem{haystacks} T.-H. Wu, G. Biamby, J. Quenum, R. Gupta, J. E. Gonzalez, T. Darrell and D. M. Chan. \href{https://arxiv.org/abs/2407.13766}{Visual Haystacks: A Vision-Centric Needle-In-A-Haystack Benchmark}. ICLR, 2025.
\bibitem{readagent} K.-H. Lee, X. Chen, H. Furuta, J. Canny and I. Fischer. \href{https://proceedings.mlr.press/v235/lee24c.html}{A Human-Inspired Reading Agent with Gist Memory of Very Long Contexts}. ICML, 2024.
\bibitem{memgpt} C. Packer, S. Wooders, K. Lin, V. Fang, S. G. Patil, I. Stoica and J. E. Gonzalez. \href{https://arxiv.org/abs/2310.08560}{MemGPT: Towards LLMs as Operating Systems}. arXiv:2310.08560, 2023.
\bibitem{mobileagente} Z. Wang et al. \href{https://arxiv.org/abs/2501.11733}{Mobile-Agent-E: Self-Evolving Mobile Assistant for Complex Tasks}. arXiv:2501.11733, 2025.
\bibitem{guirise} T. Liu, C. Wang, R. Li, Y. Yu, X. He and S. Bai. \href{https://papers.nips.cc/paper_files/paper/2025/hash/f144ab9985c739a5091ec188a2688644-Abstract-Conference.html}{GUI-Rise: Structured Reasoning and History Summarization for GUI Navigation}. NeurIPS, 2025.
\bibitem{palui} Z. Liu, J. Li, W. X. Zhao, D. Gao, Y. Li and J.-R. Wen. \href{https://aclanthology.org/2026.findings-aacl.13/}{PAL-UI: Planning with Active Look-back for Vision-Based GUI Agents}. Findings of IJCNLP-AACL, 2026.
\bibitem{qwen} Qwen Team. \href{https://arxiv.org/abs/2412.15115}{Qwen2.5 Technical Report}. arXiv:2412.15115, 2024.
\bibitem{lora} E. J. Hu, Y. Shen, P. Wallis, Z. Allen-Zhu, Y. Li, S. Wang, L. Wang and W. Chen. \href{https://openreview.net/pdf?id=nZeVKeeFYf9}{LoRA: Low-Rank Adaptation of Large Language Models}. ICLR, 2022.
\bibitem{odysseyv1} Q. Lu et al. \href{https://arxiv.org/abs/2406.08451v1}{GUI Odyssey: A Comprehensive Dataset for Cross-App GUI Navigation on Mobile Devices}. arXiv:2406.08451v1, 2024. Initial dataset release.
\bibitem{aitz} J. Zhang et al. \href{https://aclanthology.org/2024.findings-emnlp.702/}{Android in the Zoo: Chain-of-Action-Thought for GUI Agents}. Findings of EMNLP, 2024.
\bibitem{amex} Y. Chai et al. \href{https://aclanthology.org/2025.findings-acl.110/}{AMEX: Android Multi-annotation Expo Dataset for Mobile GUI Agents}. Findings of ACL, 2025.
\bibitem{androidworld} C. Rawles et al. \href{https://proceedings.iclr.cc/paper_files/paper/2025/hash/01a83bc2f2732a58e6aa731e659e7101-Abstract-Conference.html}{AndroidWorld: A Dynamic Benchmarking Environment for Autonomous Agents}. ICLR, 2025.
\bibitem{autogui} Z. Zhang and A. Zhang. \href{https://aclanthology.org/2024.findings-acl.186/}{You Only Look at Screens: Multimodal Chain-of-Action Agents}. Findings of ACL, 2024.
\bibitem{cogagent} W. Hong et al. \href{https://openaccess.thecvf.com/content/CVPR2024/html/Hong_CogAgent_A_Visual_Language_Model_for_GUI_Agents_CVPR_2024_paper.html}{CogAgent: A Visual Language Model for GUI Agents}. CVPR, 2024.
\bibitem{seeclick} K. Cheng et al. \href{https://aclanthology.org/2024.acl-long.505/}{SeeClick: Harnessing GUI Grounding for Advanced Visual GUI Agents}. ACL, 2024.
\bibitem{osatlas} Z. Wu et al. \href{https://proceedings.iclr.cc/paper_files/paper/2025/hash/0faa4bc5f522076947a030273629d4fe-Abstract-Conference.html}{OS-ATLAS: Foundation Action Model for Generalist GUI Agents}. ICLR, 2025.
\bibitem{synapse} L. Zheng, R. Wang, X. Wang and B. An. \href{https://proceedings.iclr.cc/paper_files/paper/2024/hash/52f050499cf82fa8efb588e263f6f3a7-Abstract-Conference.html}{Synapse: Trajectory-as-Exemplar Prompting with Memory for Computer Control}. ICLR, 2024.
\bibitem{awm} Z. Z. Wang, J. Mao, D. Fried and G. Neubig. \href{https://proceedings.mlr.press/v267/wang25bx.html}{Agent Workflow Memory}. ICML, 2025.
\bibitem{mintz} M. Mintz, S. Bills, R. Snow and D. Jurafsky. \href{https://aclanthology.org/P09-1113/}{Distant Supervision for Relation Extraction without Labeled Data}. ACL-IJCNLP, 2009.
\bibitem{drqa} D. Chen, A. Fisch, J. Weston and A. Bordes. \href{https://aclanthology.org/P17-1171/}{Reading Wikipedia to Answer Open-Domain Questions}. ACL, 2017.
\bibitem{orqa} K. Lee, M.-W. Chang and K. Toutanova. \href{https://aclanthology.org/P19-1612/}{Latent Retrieval for Weakly Supervised Open Domain Question Answering}. ACL, 2019.
\bibitem{dpr} V. Karpukhin, B. Oguz, S. Min, P. Lewis, L. Wu, S. Edunov, D. Chen and W.-t. Yih. \href{https://aclanthology.org/2020.emnlp-main.550/}{Dense Passage Retrieval for Open-Domain Question Answering}. EMNLP, 2020.
\bibitem{monot5} R. Nogueira, Z. Jiang, R. Pradeep and J. Lin. \href{https://aclanthology.org/2020.findings-emnlp.63/}{Document Ranking with a Pretrained Sequence-to-Sequence Model}. Findings of EMNLP, 2020.
\bibitem{rankgpt} W. Sun et al. \href{https://aclanthology.org/2023.emnlp-main.923/}{Is ChatGPT Good at Search? Investigating Large Language Models as Re-Ranking Agents}. EMNLP, 2023.
\end{thebibliography}
\endgroup
